\section{线性变换与换基：对象不变，坐标可以改变}
\label{sec:03-Linear-Algebra/02-Vector-Spaces/04}

前面我们学会了用基把一个向量写成一组数字。接下来的麻烦是：如果一个操作把一个向量变成另一个向量，用不同的基写它的矩阵，得到的数字会完全不一样。那哪一列数字才是这个操作真正的面目？

都不是。矩阵只是操作在特定坐标系下的照片。换个角度再拍，照片里的数字会变，但被拍的操作还是同一个。这节就是要把操作和它的照片分开，然后说明换照片的角度不是为了制造混乱，而是为了让操作的骨架更清楚。

\subsection{线性变换：两个规则定义一类最规矩的操作}

映射 \(T:V\to W\) 被称为线性变换，满足
\[
T(\alpha u+\beta v)=\alpha T(u)+\beta T(v)
\]
对所有 \(u,v\in V\) 和标量 \(\alpha,\beta\) 成立。两条规则拆开看：缩放能穿过 \(T\)（齐次性），加法能穿过 \(T\)（可加性）。合在一起等价于一个更省的表述：\(T\) 保持所有线性组合。

直接推出 \(T(0)=0\)。所以 \(T(x)=Ax+b\) 当 \(b\neq0\) 时不是线性变换——保持原点是线性性的硬门槛。\hyperref[sec:03-Linear-Algebra/02-Vector-Spaces/01]{``为什么需要向量空间''}一节里的非齐次方程解集 \(x_p+\mathcal N(A)\)，就是子空间的平移，正是跨过了原点这条线。

矩阵乘法当然是线性变换，但不是全部。下面这些同样是线性变换：
\begin{itemize}
\tightlist
\item
  多项式求导：\(D(p)=p'\)，满足叠加和缩放的穿过；
\item
  定积分：固定上下限的积分是线性泛函；
\item
  卷积：固定核的卷积算子；
\item
  转置：把矩阵映射到其转置；
\item
  投影到子空间：沿正交方向压缩维度。
\end{itemize}

矩阵之所以在所有地方冒出来，是因为有限维空间中，任何线性变换一旦选定了输入基和输出基，就可以写成一个矩阵。矩阵是线性变换的有限维坐标表示——它不是变换本身，却可以替变换完成全部计算。

\subsection{知道基向量去了哪里，就知道了一切}

这是线性变换最省力的性质。设 \(B=(v_1,\ldots,v_n)\) 是 \(V\) 的一组基，任意 \(v\in V\) 唯一写成
\[
v=c_1v_1+\cdots+c_nv_n.
\]
线性性直接给出：
\[
T(v)=c_1T(v_1)+\cdots+c_nT(v_n).
\]
只要对每个基向量运行一次 \(T\)，记录下结果，整个变换在全体输入上的行为就全部被锁定了。有限维空间把查看无限多点的问题缩减为查看有限个基向量——这就是维数作为自由度上限的另一个面孔。

在此基础上，再为输出空间 \(W\) 选一组基 \(C\)，把每个 \(T(v_j)\) 的 \(C\) 坐标填入矩阵的第 \(j\) 列：
\[
[T]_{C\leftarrow B}.
\]
这个矩阵满足
\[
[T(v)]_C=[T]_{C\leftarrow B}[v]_B.
\]
箭头方向故意写成从 B 指向 C：矩阵吃进 B 坐标，吐出 C 坐标。记住箭头的朝向，就能避免换基时左侧该乘还是右侧该乘这种常规混淆。

\subsection{求导算子的矩阵：不是方阵也可以}

取 \(D:P_2\to P_1\)，\(D(p)=p'\)。输入基
\[
B=(1,t,t^2),
\]
输出基
\[
C=(1,t).
\]
依次求导：
\[
D(1)=0,\quad D(t)=1,\quad D(t^2)=2t.
\]
三个基向量的像是 \(0\)（零多项式）、\(1\)、\(2t\)。把它们的 \(C\) 坐标作为列：
\[
[D]_{C\leftarrow B}=
\begin{pmatrix}
0&1&0\\
0&0&2
\end{pmatrix}.
\]
检查这个矩阵是否做对了事：若 \(p(t)=a+bt+ct^2\)，其 \(B\) 坐标为 \((a,b,c)^\trans\)。矩阵乘法：
\[
\begin{pmatrix}0&1&0\\0&0&2\end{pmatrix}
\begin{pmatrix}a\\b\\c\end{pmatrix}=
\begin{pmatrix}b\\2c\end{pmatrix},
\]
对应的多项式恰好是 \(b\cdot1+2c\cdot t=p'(t)\)。

矩阵不是方阵——\(2\times3\)，因为输入空间是 3 维，输出空间是 2 维。维数的变化直接刻在矩阵的尺寸上。

\subsection{核与像：线性变换自带的账目}

核和像推广了矩阵的零空间和列空间：
\[
\ker T=\{v\in V:T(v)=0\},\qquad
\operatorname{im}T=\{T(v):v\in V\}.
\]
核在输入端，记录被消灭的部分；像在输出端，记录被产生的部分。

在求导的例子中，常数函数的导数是零，所以
\[
\ker D=\operatorname{span}\{1\},
\]
维数为 1。每个一次多项式 \(u+vt\) 都可以通过积分找回原像：
\[
D\!\left(ut+\frac{v}{2}t^2\right)=u+vt,
\]
因此
\[
\operatorname{im}D=P_1,
\]
维数为 2。维数账目闭合：
\[
\dim P_2=\dim\ker D+\dim\operatorname{im}D=1+2=3.
\]

核只含零向量时，变换不会把两个不同输入合并——这是单射。像填满整个输出空间时，每个输出都可到达——这是满射。如果输入输出维数相同，秩--零度定理带来的结论很漂亮：单射和满射等价，两者又等价于表示矩阵可逆。因为 \(\dim\ker T=0\) 与 \(\dim\operatorname{im}T=\dim W\) 在同维数设定下要么同时成立，要么同时失败。

\subsection{坐标变换：你以为换了空间，其实只是换了测量单位}

设 \(B=(v_1,\ldots,v_n)\) 是 \(\R^n\) 的一组基。把这些基向量的标准坐标排成矩阵：
\[
S=\begin{pmatrix}v_1&\cdots&v_n\end{pmatrix}.
\]
\(S\) 的角色就是翻译器：把新基坐标翻成标准坐标。若 \([x]_B\) 是 \(x\) 的 \(B\) 坐标，
\[
x=S[x]_B.
\]
反过来，
\[
[x]_B=S^{-1}x.
\]
实际操作中最容易出错的地方不是公式本身，而是搞混当前看到的数字属于哪套坐标。一个实用的习惯：在每一步计算旁标注这行数字是哪组基下的坐标。方向自然跟着标注走。

拿一个具体例子摸一遍。取
\[
v_1=\begin{pmatrix}1\\1\end{pmatrix},
\qquad
v_2=\begin{pmatrix}1\\-1\end{pmatrix},
\]
则
\[
S=\begin{pmatrix}1&1\\1&-1\end{pmatrix}.
\]
若 \([x]_B=(3,1)^\trans\)，标准坐标就是
\[
x=S[x]_B=3v_1+v_2=\begin{pmatrix}4\\2\end{pmatrix}.
\]

用路网来想：\(v_1\) 指向东北，走一段意味着向东一个单位、向北一个单位；\(v_2\) 指向东南，走一段意味着向东一个单位、向北负一个单位（也就是向南一个单位）。\([x]_B=(3,1)^\trans\) 表示沿 \(v_1\) 走三段、沿 \(v_2\) 走一段。换算成地图上的标准坐标：向东 \(3+1=4\)，向北 \(3-1=2\)。同一个位移，两套坐标系里的数字完全不同。

这里有一个值得注意的陷阱：如果 \(v_1,v_2\) 不是正交归一的（事实上这里 \(v_1^\trans v_2=0\)，它们是正交的但长度都是 \(\sqrt{2}\)），那么坐标数值本身不能直接当欧氏距离用。\"走了三格\"的长度取决于基向量的长度。后文正交单元会专门处理这个问题。

\subsection{同一变换为何对应不同的矩阵}

设 \(T:\R^n\to\R^m\) 在标准基下的矩阵为 \(A\)。现在给输入端换基 \(B\)（基矩阵 \(S\)），输出端换基 \(C\)（基矩阵 \(R\)）。从输入的 \(B\) 坐标出发，走一遍完整流程：

\[
[x]_B\xrightarrow{S}x\xrightarrow{A}T(x)\xrightarrow{R^{-1}}[T(x)]_C.
\]
先翻译到标准坐标，执行变换，再翻译到新输出坐标。因此新表示矩阵是
\[
[T]_{C\leftarrow B}=R^{-1}AS.
\]
三个因子各自只做一件事：\(S\) 翻译坐标，\(A\) 执行变换，\(R^{-1}\) 再次翻译。

两端都回到同一个空间且使用同一组新基时，
\[
A_B=S^{-1}AS.
\]
这叫相似变换。\(A\) 和 \(A_B\) 数字可以截然不同，却表示同一个线性作用，因此共享所有不依赖坐标的性质：秩、行列式、特征值、迹——后文各单元会逐一遇到。

\subsection{好基会使什么变容易}

换基不是数字游戏。它的全部动机在于：用合适的基，让原本绞在一起的结构散开。

如果能找到特征向量作为基，矩阵变成对角形——变换变成每个方向独立缩放；使用正交归一基，坐标不需要解方程组，直接点积就能拿到；在 SVD 中，输入侧和输出侧分别选各自最好的正交基后，任意矩阵的行为简化为一组独立缩放通道。

举个小例子把抽象拉回地面。一张彩色像素的标准坐标是 \((R,G,B)\)，分别记录红绿蓝三个通道的强度。换个基试试：改用
\[
\begin{pmatrix}1\\1\\1\end{pmatrix},\;
\begin{pmatrix}1\\-1\\0\end{pmatrix},\;
\begin{pmatrix}1\\1\\-2\end{pmatrix}.
\]
第一个方向是总亮度，后两个方向只改变颜色对比不改变整体亮度。同一个像素在这组新基下的坐标变成了（亮度，红偏绿，蓝偏其他）。这就是 JPEG 色彩空间转换背后的几何直觉——把整体亮度和颜色偏移拆开，然后对不同方向用不同精度压缩。实际色彩系统会使用经过感知校准的系数矩阵，但几何上就是一次换基。

适合问题结构的基让你只盯住变化的方向，而不是在每个分量上同时费力。所以换基与其说是计算技巧，不如说是观察视角的选择：你选择看变换的哪一面。

另一个快例在多项式求导。\(D:P_3\to P_3\) 在幂基 \((1,t,t^2,t^3)\) 下的矩阵是
\[
\begin{pmatrix}
0&1&0&0\\
0&0&2&0\\
0&0&0&3\\
0&0&0&0
\end{pmatrix}.
\]
这个矩阵看起来没什么规律——每列的数字在往右上角漂。但换到 Hermite 多项式构成的基之后（在函数空间的标准内积下是正交基），同样一个求导算子的矩阵会变成三对角且结构规律。矩阵的形式取决于选什么样的基去观察同一个操作。

更有意思的是，合适基下矩阵的稀疏性、对角优势等结构性质直接决定了数值算法的稳定性和迭代速度。大规模线性系统中，很少人直接用随手拿到的标准基矩阵做计算——都会先找一个使矩阵的谱性质更有利的基，这本身就是一次换基。所以换基这笔账早晚要算，区别只在于主动还是被动。

\subsection{变换与表示之间隔着什么}

线性变换是空间之间的操作。矩阵是这个操作在选定坐标系下的数字记账。两者之间的关系就像物体和照片——同一物体可以进行不同角度的拍摄，得到完全不同的二维投影。但物体不会因为换了一个角度就变成另一个东西。

这个区分对初学有两层实用提醒。第一，看到两个数字不同的矩阵，不要立刻断定它们代表不同的变换——它们可能是同一变换在不同基下的不同记账方式。相似变换 \(S^{-1}AS\) 就专做这类识别。第二，看到矩阵某个位置为零不要急于下结论——换一组基之后那个位置可能就不再是零了。矩阵的零元素既可能反映变换的结构，也可能只是坐标系带来的偶然。

整个线性代数的后半程——特征值分解、奇异值分解、Jordan 标准形——全都在做同一件事：为一个线性变换寻找使其矩阵最简单的那组基。坐标可以用不同的数字去记录同一个操作，但这些数字里藏着简化它的钥匙。最有用的那组基，往往就是让矩阵非零元最少、结构最规律的那组。

延伸阅读：MIT 18.06：Linear Transformations and their Matrices。
